\documentclass[11pt,a4paper]{article}
\usepackage[utf8]{inputenc}
\usepackage[T1]{fontenc}
\usepackage[brazilian]{babel}
\usepackage{geometry}
\usepackage{graphicx}
\usepackage{booktabs}
\usepackage{float}
\usepackage{enumitem}
\usepackage{amsmath}
\usepackage{xcolor}
\usepackage{tcolorbox}
\geometry{margin=2.5cm}
\setlength{\parindent}{0pt}
\setlength{\parskip}{0.5em}
\tcbuselibrary{skins,breakable}
\newtcolorbox{nota}[1][]{colback=yellow!10,colframe=orange!60!black,fonttitle=\bfseries,title=#1,breakable}

\begin{document}
\begin{center}
  
\includegraphics[width=0.16\textwidth]{../../../figures_global/unb.png}\\[0.3cm]
  {\large\textbf{GABARITO DO PROFESSOR -- USO RESTRITO}}\\[0.2cm]
  {\large FGA0092 -- Princípios de Comunicação para Engenharia}\\[0.2cm]
  \rule{12cm}{0.8pt}\\[0.2cm]
  {\Large\textbf{Prática 08: Amostragem e Aliasing}}\\[0.2cm]
  \rule{12cm}{0.8pt}
\end{center}

\section*{Gráfico esperado}
\begin{center}
  \includegraphics[width=0.88\textwidth]{grafico_esperado.png}
\end{center}
\textit{O gráfico acima foi gerado pelo script \texttt{gerar\_grafico.py} desta pasta.}

\section*{Respostas esperadas}

\begin{enumerate}[label=\textbf{Q\arabic*.}, leftmargin=*, itemsep=1em]

  \item \textbf{Verifique o critério de Nyquist e calcule a frequência de aliasing para os cenários A--F.}

    \textbf{Resposta esperada:}
    Cenário A ($f_0=1000$, $f_{s,ef}=12000$\,Hz): Nyquist ok ($1000 < 6000$); sem aliasing.
    Cenário B ($f_0=5000$, $f_{s,ef}=12000$\,Hz): Nyquist ok ($5000 < 6000$); sem aliasing.
    Cenário C ($f_0=7000$, $f_{s,ef}=12000$\,Hz): violação ($7000 > 6000$);
    $f_\text{alias} = |7000 - 12000| = 5000$\,Hz.
    Cenário D ($f_0=1000$, $f_{s,ef}=4800$\,Hz): ok ($1000 < 2400$); sem aliasing.
    Cenário E ($f_0=2500$, $f_{s,ef}=4800$\,Hz): violação ($2500 > 2400$);
    $f_\text{alias} = |2500 - 4800| = 2300$\,Hz.
    Cenário F ($f_0=3500$, $f_{s,ef}=4800$\,Hz): violação;
    $f_\text{alias} = |3500 - 4800| = 1300$\,Hz.

    \begin{nota}[Erros comuns]
      Calcular $f_\text{alias}$ como $f_0 - f_s$ em vez de $|f_0 - n f_s|$ para o
      valor inteiro $n$ que minimiza o resultado, ou usar $f_s$ em vez de $f_s/2$
      como frequência de Nyquist.
    \end{nota}

  \item \textbf{Como o aliasing se manifesta no domínio do tempo no Cenário C?}

    \textbf{Resposta esperada:} No Cenário C, o sinal de 7000\,Hz após decimação
    aparenta ter frequência 5000\,Hz. No tempo, a forma de onda oscila mais lentamente
    que o sinal original: em vez de 7 ciclos por milissegundo, parece ter 5 ciclos.
    Isso ocorre porque as amostras capturadas a 12000\,Hz de um sinal de 7000\,Hz
    correspondem exatamente às amostras que seriam obtidas de um sinal de 5000\,Hz
    (reflexo em $f_s/2 = 6000$\,Hz).

    \begin{nota}[Erros comuns]
      Descrever o aliasing apenas no espectro sem relacionar com a manifestação
      no domínio do tempo: o sinal parece oscilar em frequência diferente da real.
    \end{nota}

  \item \textbf{Qual é a diferença entre frequência real e frequência aparente após amostragem insuficiente?}

    \textbf{Resposta esperada:} A frequência aparente é a frequência com que as amostras
    parecem variar; a frequência real é a do sinal analógico original. Sem informação
    adicional (taxa de amostragem e frequência do sinal fonte), um receptor não pode
    distinguir 7000\,Hz subamostrado de um sinal real de 5000\,Hz: ambos produzem
    exatamente as mesmas sequências de amostras. Esta ambiguidade é a essência do aliasing
    e justifica o filtro anti-aliasing antes da amostragem.

    \begin{nota}[Erros comuns]
      Dizer que o receptor ``pode calcular'' a frequência real sem informação adicional.
      A ambiguidade é fundamental e não pode ser resolvida apenas com as amostras.
    \end{nota}

  \item \textbf{Como o filtro anti-aliasing (LPF) resolve o problema nos Cenários C e E? Há alguma limitação?}

    \textbf{Resposta esperada:} No Cenário C: o LPF com \texttt{cutoff}$=6000$\,Hz atenua
    o sinal de 7000\,Hz antes da decimação, eliminando o aliasing; o pico em 5000\,Hz
    desaparece. No Cenário E: o LPF com \texttt{cutoff}$=2400$\,Hz deve atenuar o sinal
    de 2500\,Hz. Como $f_0=2500 > f_\text{cutoff}=2400$\,Hz, o próprio sinal de interesse
    é atenuado -- o filtro AA resolve o aliasing mas remove o sinal que está acima de
    $f_s/2$, incluindo o sinal útil quando ele está muito próximo da frequência de Nyquist.

    \begin{nota}[Erros comuns]
      Afirmar que o LPF ``preserva o sinal'' no Cenário E sem notar que
      $2500 > 2400$\,Hz e o sinal de interesse é atenuado pelo próprio filtro AA.
    \end{nota}

  \item \textbf{Quais são as diferenças entre o bloco \texttt{Keep 1 in N} (simulação) e um ADC real?}

    \textbf{Resposta esperada:} O bloco \texttt{Keep 1 in N} realiza subamostragem ideal:
    descarta amostras sem ruído de quantização, sem jitter de amostragem, sem distorção de
    abertura. Um ADC real: (1) tem jitter de relógio que limita o SNR
    ($\text{SNR}_\text{jitter} = 1/(2\pi f_0 \tau_j)^2$); (2) tem largura de abertura finita;
    (3) tem ruído de quantização proporcional ao número de bits; (4) tem largura de banda
    de entrada limitada pelo circuito amostrador. Na simulação, o aliasing é perfeitamente
    determinístico; no ADC real, não-idealidades adicionam aleatoriedade.

    \begin{nota}[Erros comuns]
      Descrever apenas o jitter sem mencionar o ruído de quantização ou a largura de
      banda de entrada, que são também limitações fundamentais do ADC real.
    \end{nota}

\end{enumerate}

\section*{Critérios de correção}
\begin{tabular}{ll}
  \toprule
  Critério & Pontuação \\
  \midrule
  Q1 -- Critério de Nyquist e cálculo de $f_\text{alias}$ para todos os cenários & 1,0 ponto \\
  Q2 -- Manifestação do aliasing no tempo: oscilação mais lenta, não só espectro  & 1,0 ponto \\
  Q3 -- Ambiguidade fundamental: amostras idênticas para duas frequências distintas & 1,0 ponto \\
  Q4 -- Limitação do filtro AA no Cenário E: atenua o próprio sinal de interesse  & 1,0 ponto \\
  Q5 -- Diferenças entre Keep 1 in N e ADC real: jitter, quantização e banda      & 1,0 ponto \\
  \midrule
  \textbf{Total} & \textbf{5,0 pontos} \\
  \bottomrule
\end{tabular}

\bigskip
\begin{nota}[Nota de correção geral]
  Priorizar coerência física e evidência experimental (capturas de tela, tabelas).
  Respostas sem evidência visual devem ter desconto. Não penalizar pequenos erros
  de cálculo se o raciocínio físico estiver correto.
\end{nota}

\end{document}
